\documentclass[12pt]{article}
\usepackage[utf8]{inputenc}
\usepackage[T1]{fontenc}
\usepackage{hyperref}
\title{Ani Phowa Language Resources for K-12 Curriculum Development}
\author{Compiled by Claude Sonnet 5.5 (claude-sonnet-5-5) for Andr\'as Kornai}
\date{}
\begin{document}
\maketitle

\section{Introduction}
Ani Phowa (ISO 639-3 ypn; Glottocode anip1235) is a Loloish (Sino-Tibetan) language. Asia Harvest places its speakers in Yanggai District, Kaiyuan County, Honghe Prefecture, Yunnan. Wikipedia lists Ani as one of three Phowa varieties (Helpho, Ani, Labo) and gives about 10,000 speakers. Joshua Project gives 9,900. The Ani are officially counted within the Yi nationality. Asia Harvest reports rapid assimilation to Han language and culture, and Wikipedia says use is decreasing.

\section{Primary Educational Resources}
No primers, textbooks or school programmes for Ani Phowa turned up in the searches.

\section{Linguistic and Reference Materials}
\begin{itemize}
\item Wikipedia, Phowa language: \url{https://en.wikipedia.org/wiki/Phowa_language}
\item Wikipedia, Loloish languages: \url{https://en.wikipedia.org/wiki/Loloish_languages}
\item Wikipedia, Southeastern Loloish languages (found under a mirror URL): \url{https://www.Wikipedia.com/wiki/Phula_language}
\item Pelkey, ``Puzzling over Phula'', on Academia.edu (found in results; I did not check whether it covers Ani Phowa in detail): \url{https://www.academia.edu/2435607/Puzzling_over_Phula}
\end{itemize}

\section{Religious and Community Resources}
\begin{itemize}
\item Asia Harvest people-group profile of the Ani in Yunnan (PDF): \url{https://people-groups.asiaharvest.org/people-group-profiles/yunnan/ani.pdf}
\item Joshua Project, Ani Phowa in China: \url{https://joshuaproject.net/people_groups/20627}
\item Joshua Project, Ani in China. One Joshua Project profile places the group in Vietnam as well, which I could not verify: \url{https://joshuaproject.net/people_groups/18396}
\end{itemize}
Both organisations are missionary-oriented. No scripture translation was confirmed.

\section{Digital and Audio Resources}
No audio, video or corpus resources were found.

\section{Curriculum Development Context}
Very little is available. The language appears to be shrinking in the face of Han assimilation, which makes documentation and community-led work urgent. Materials for related Southeastern Loloish varieties may help as a starting point but should not be assumed to transfer.
\end{document}